% \subsubsection*{Ethics Statement}
% The authors have read and adhered to the ICLR Code of Ethics. Our work introduces a powerful video generation model, EchoMotion, and a new human-centric dataset, \textit{HuMoVe}, which prompts careful consideration of several ethical dimensions.

% \paragraph{Potential for Misuse} We acknowledge that generative models capable of creating realistic human videos, such as EchoMotion, have dual-use potential. This technology could be misused for creating synthetic media for malicious purposes (e.g., "deepfakes," disinformation). Our primary research goal is to advance the scientific understanding of motion dynamics in generative models, with potential positive applications in creative industries, virtual reality, and data augmentation for training robust perception models. We are committed to responsible research dissemination and will explore methods like digital watermarking in future work to mitigate misuse.

% \paragraph{Bias in Data and Models} The \textit{HuMoVe} dataset, despite our efforts to capture diverse motions, may reflect biases present in the source data regarding demographics, cultural representation, and types of activities. Consequently, models trained on this data, including EchoMotion, may inherit and potentially amplify these biases. We recognize this as a limitation and an important direction for future work, which should focus on dataset expansion, bias detection, and algorithmic fairness to ensure more equitable generative outcomes.


% \subsection*{Reproducibility Statement}
% We are committed to ensuring the reproducibility of our research. To this end, we provide detailed information and resources as described below.

% \paragraph{Code and Pre-trained Models} The complete source code for our EchoMotion model, along with the pre-trained model weights, will be publicly released in a GitHub repository upon publication. This will allow the community to verify our results, build upon our work, and apply the model to new tasks.

% \paragraph{Dataset} Our newly constructed \textit{HuMoVe} dataset will be publicly released. This release will include the video identifiers, our detailed textual annotations, and the extracted SMPL parameters. We provide a comprehensive description of the data collection, filtering, and annotation process in the Appendix \ref{build_data} to ensure transparency.


% \paragraph{Computational Environment} The Appendix includes a description of the computational environment used for our experiments, specifying key dependencies such as PyTorch version, CUDA version, and other relevant libraries, to help other researchers replicate our setup.